\documentclass{ieeeaccess}
\usepackage{cite}
\usepackage{amsmath,amssymb,amsfonts}
\usepackage{graphicx}
\usepackage{textcomp}
\usepackage{booktabs}
\usepackage{multirow}
\usepackage{adjustbox}
\usepackage[hidelinks]{hyperref}

\graphicspath{{../../figures/ml/}{../../figures/}}

% Headline numbers (test, TH + VN, Delta = 30, unconditional protocol). Every value comes from docs/report/tables/ml/;
% update this block and the tables below when a rerun replaces a table.
\newcommand{\nRefCond}{0.206}        % t28, P1-frozen, conditional NDCG@20
\newcommand{\nRefUncond}{0.095}      % t28, P1-frozen, unconditional NDCG@20, TH + VN
\newcommand{\nNoPosShare}{40.7}      % t28, % of anchors with no positive among the 200 candidates
\newcommand{\nNoSource}{58}          % t28b, % of positives proposed by no source
\newcommand{\nCandOld}{0.370}        % t30, R1-R5@200 coverage
\newcommand{\nCandNew}{0.450}        % t30, candidates_v2 coverage
\newcommand{\nExpOld}{0.307}         % t30, R1-R5@200 explore coverage
\newcommand{\nExpNew}{0.400}         % t30, candidates_v2 explore coverage
\newcommand{\nK}{300}                % candidate cap of candidates_v2
\newcommand{\nBest}{0.1020}          % t31, best model on valid (P3 one head)
\newcommand{\nTuned}{0.1011}         % t31, P1-tuned
\newcommand{\nBaseline}{0.0788}      % t31, best next-basket baseline (TIFU-KNN)

\def\footerline{IT2043 Advanced Machine Learning and Applications, final report, 2026}

\begin{document}
\history{Final report of the course Advanced Machine Learning and Applications (IT2043.26.1.CH.01), Master of Information Technology,
November 2026. Source code: github.com/haiiyins/retail\_multidim\_pipeline}
\doi{}

\title{Evaluating and Improving Personalized Product Recommendation with Multi-Signal Machine Learning on Large-Scale Retail
Data}

\author{\uppercase{Yen Phan}\authorrefmark{1,*}, \uppercase{Manh Le}\authorrefmark{1,*}, and
\uppercase{Kha Nguyen}\authorrefmark{1,*}}
\address[1]{University of Information Technology, Vietnam National University, Ho Chi Minh City, Vietnam}
\address[*]{Corresponding authors}

\markboth
{Y. Phan \headeretal: Evaluating and Improving Personalized Product Recommendation}
{Y. Phan \headeretal: Evaluating and Improving Personalized Product Recommendation}

\corresp{Corresponding authors: Yen Phan (e-mail: 26848274@grad.uit.edu.vn), Manh Le (e-mail: 26848126@grad.uit.edu.vn), and
Kha Nguyen (e-mail: 26848091@grad.uit.edu.vn).}

\begin{abstract}
An omnichannel retailer records purchases, web behaviour, prices, promotions, stock and stores, yet recommends the same
best-sellers or past purchases to everyone. On eleven operational exports of a multi-country sporting-goods retailer (66.1\,M sale
lines, 1.0\,B web events, 4.7\,M buying members), a Bronze / Silver / Gold lakehouse with seven quality gates builds a point-in-time
feature store with zero leakage violations. Scored, as usual, only on members whose candidates hold a purchase, a rule-based
two-stage recommender reaches an NDCG@20 of \nRefCond{}; counting every purchase it reaches \nRefUncond{}, and \nNoSource{}\,\% of the
test purchases were never retrieved, almost all items new to the member. We therefore treat retrieval as a model: eight learned
sources restricted to new items, refitted quarterly and mixed under a cap of \nK{} candidates, raise coverage from \nCandOld{} to
\nCandNew{}. Under a uniform 10-trial tuning budget, gradient-boosted, stacked and deep multi-task rankers all beat five next-basket
baseline families by more than a quarter; features give about 70\,\% of the gain over the baseline ranker, tuning the rest, and the
model family least. A single network over the whole catalogue matches the two-stage pipeline, and its rank fusion with the
pipeline, pooled over three countries, reaches 0.1102; predicting the product type instead of the model lifts NDCG@20 to 0.372. On
purchase-log features alone the pipeline keeps about four fifths of its gain over frequency rules, and naive stacking inflates
in-sample metrics.
\end{abstract}

\begin{keywords}
Data lakehouse, evaluation methodology, gradient boosting, learning to rank, multi-task learning, next-basket recommendation, recommender systems,
retail analytics, two-stage retrieval.
\end{keywords}

\titlepgskip=-15pt

\maketitle

% ---------------------------------------------------------------------------------------------------------------------------------
\section{Introduction}
\label{sec:introduction}

\PARstart{A}{n} omnichannel retailer knows a great deal about its loyalty members: what they bought and returned, what they looked at
on the web site and in the app, which sports they follow, and, for every product, its price, promotions, stock and season. Yet the
list a member sees is usually the same for everyone, the national best-sellers, or a reminder to buy the same items again. The
business goal of this work is a personal top-$K$ list for every member, every week, and an honest number for how good it is.

Three obstacles stand between the raw data and that list. \emph{Data}: each signal lives in its own operational table with its own
grain, keys and clock, and the click-stream alone holds a billion events. \emph{Time}: a feature may only use what was known before
the recommendation was made, otherwise offline scores promise gains that vanish in production. \emph{Measurement}: a score computed
only where the candidate list already holds a purchase says nothing about the members the list misses.

\textit{Problem.} Given a member $u$ at an anchor $t$ (a Monday), predict the product models $u$ buys in $(t, t+30]$ as an ordered
top-$K$ list; a hit is a model bought and not returned, and every feature uses only data recorded before $t$. A natural baseline is a
rule-based two-stage recommender: five hand-written candidate rules (R1 repeat, R2 co-purchase, R3 web department, R4 sport $\times$
offer, R5 national best-sellers) propose 200 models per (member, Monday) anchor, and a LightGBM LambdaRank
ranker~\cite{ke2017lightgbm,burges2010lambdamart} (P1), tuned once over 73 Gold pair features, orders them. Scored the usual way, on
the anchors whose list holds a purchase, it reaches a normalized discounted cumulative gain NDCG@20 of \nRefCond{}. Scored over every
anchor with a purchase, counting purchases outside the list as misses, the same model scores 0.090 (TH + VN: \nRefUncond{}), and
\nNoPosShare{}\,\% of the anchors have no purchase in the list at all. That gap is the research problem.

This paper therefore builds the data foundation first and then measures what each modelling step is worth. It turns eleven
operational exports into a leakage-controlled lakehouse, reformulates next-basket recommendation (NBR) as a two-stage problem in
which retrieval is a model with its own objective, and compares gradient-boosted trees (GBDT), a stacked variant and a deep
multi-task ranker with strong NBR baselines under one candidate list, one protocol and one tuning budget. The research questions are
phrased as questions, not as expected orderings, because a GBDT win is a valid outcome~\cite{grinsztajn2022tabular}:
\begin{itemize}
\item \textbf{RQ1 (data).} Can eleven raw operational exports become a point-in-time, leakage-free feature store at member--week
grain on one machine, and what does it hand to the models?
\item \textbf{RQ2 (retrieval).} Scored on every purchase, where does the rule-based baseline fail, and how much do learned candidate
sources raise candidate recall and explore coverage over R1--R5, at what cost, and how much of that gain reaches NDCG@20 when the
ranker is unchanged?
\item \textbf{RQ3 (ranking and tuning).} On the same candidates and the same budget, which ranker wins (tuned GBDT, stacking, a deep
multi-task ranker with repeat and explore heads, NBR baselines), and how much of the gain comes from features, hyper-parameter tuning
and the model family?
\item \textbf{RQ4 (value).} How many times better than ``best-sellers'' and ``buy the same again'' is the best system, for whom, at
what compute cost, what purchase value does it capture (offline proxy), and how much of its gain needs data beyond the purchase
history?
\item \textbf{H1 (hypothesis).} Naive stacking inflates metrics on training and pseudo-test anchors relative to a rolling protocol;
on the real test the gap narrows but does not vanish.
\end{itemize}

The contributions are: (1) a Bronze / Silver / Gold lakehouse with seven quality gates and a point-in-time ledger with zero leakage
violations; (2) an unconditional protocol with a repeat / explore split and an error decomposition that locates \nNoSource{}\,\% of
the test purchases in retrieval rather than ranking; (3) an explore-only retrieval layer of eight learned sources refitted
quarterly, with a rank-quota mixer; (4) a controlled comparison of rankers and five NBR baseline families that separates features,
tuning and model family, with a measurement of stacking leakage; (5) a single-stage network over the whole catalogue and its
fusion with the pipeline; and (6) a static demonstrator with privacy safeguards.

% ---------------------------------------------------------------------------------------------------------------------------------
\section{Related Work}
\label{sec:related}

\textit{Lakehouse and leakage control.} The lakehouse architecture~\cite{armbrust2021lakehouse} keeps open columnar files as the
storage layer and adds managed tables and transformations on top; DuckDB~\cite{raasveldt2019duckdb} runs such analytical SQL in
process on one machine. Leakage, a feature that uses information unavailable at prediction time, is a common cause of inflated and
irreproducible results in applied machine learning~\cite{kapoor2023leakage}; we test for it at every anchor.

\textit{Next-basket recommendation.} The reality check of Li et al.~\cite{li2023nbr} showed that repeat purchases dominate common NBR
metrics and that simple frequency baselines (P-TopFreq, GP-TopFreq) are hard to beat; it introduced the repeat / explore recall split
used here. TIFU-KNN~\cite{hu2020tifu} models personal item frequency with time decay and neighbour smoothing. More generally,
carefully tuned simple baselines often match recent neural recommenders~\cite{dacrema2019progress}.

\textit{Item-based and graph retrieval.} EASE~\cite{steck2019ease} is a closed-form linear auto-encoder; RP3beta~\cite{paudel2017rp3beta}
re-weights random walks by item popularity; item2vec~\cite{barkan2016item2vec} learns co-occurrence embeddings.

\textit{Neural retrieval and sequential models.} Two-tower models~\cite{yi2019twotower} score users and items in a shared space with a
softmax over items; SASRec~\cite{kang2018sasrec} applies causal self-attention~\cite{vaswani2017attention} to the behaviour sequence,
gSASRec~\cite{petrov2023gsasrec} corrects its overconfidence under negative sampling, and PinnerFormer~\cite{pancha2022pinnerformer}
trains on a dense all-action objective over a future window. DIN~\cite{zhou2018din} adds target attention of the candidate over the
behaviour sequence.

\textit{Multi-task and repeat-aware ranking.} MMoE~\cite{ma2018mmoe} gates shared experts per task. Repeat-aware recommenders separate
repeat and explore modes (RepeatNet~\cite{ren2019repeatnet}, ReCANet~\cite{ariannezhad2022recanet}). Our P3 gates two listwise heads
with a learned member-level probability.

\textit{Leakage in stacking.} Out-of-fold or temporally rolling fits are standard for stacked models~\cite{wolpert1992stacking}; we
measure the inflation of an in-sample fit directly.

% ---------------------------------------------------------------------------------------------------------------------------------
\section{Data Foundation}
\label{sec:data}

This section answers RQ1: how the raw exports become model inputs, and what the processing leaves for the models.
Fig.~\ref{fig:overview} gives the overview.

\Figure[t!](topskip=0pt, botskip=0pt, midskip=0pt)[width=0.95\textwidth]{architecture_overview.png}{Data foundation: Bronze, Silver and Gold tiers of the DuckDB lakehouse, with the quality gates QC-1 to QC-7 across every layer.\label{fig:overview}}

\subsection{Sources}
The data set is an export of eleven operational tables for Singapore (SG), Thailand (TH) and Vietnam (VN) from 2023-01 to 2026-10,
about 230\,GB raw: transactions (89.3\,M sale, cancel and return lines, 29.7\,M payment lines, 63.0\,M transaction--member links),
master data (5.7\,M loyalty accounts, 4.4\,M SKUs, 275 stores and web sites, a calendar), price and stock timelines (42.4\,M and
36.7\,M rows), 0.8\,M favourite-sport events, and 999.7\,M web and app events in 1\,774 parts. Every table has its own grain, keys and
clock. Nothing was transformed before its assumptions were measured: a full scan of every file found the facts that shaped the
processing, among them 9.5\,\% duplicate sales rows from overlapping extracts, cancellations recorded as separate lines that net a sale
to zero, sale times in local clocks while version stamps are in UTC, only 34\,\% of the web events carrying a member, a stock table
that tracks spare parts only (the stock signal is therefore switched off), and store fixtures and carrier bags in the catalogue.

\subsection{Through the Tiers}
\textit{Bronze keeps.} A data contract per table declares its key, required columns, join-key types and partition. The loader writes
ZSTD Parquet partitioned by country, year and month (127\,GB, full load in 1\,h\,51\,min, the click-stream alone 87\,min); rows the
parser cannot read go to a quarantine table with a reason instead of being dropped (65 rows).

\textit{Silver corrects.} Fifteen dbt models clean and harmonise, one technique per data problem: versioned deduplication and signed
sale $-$ cancel netting per transaction and item, with the status derived from the net amount; UTC through each unit's time zone,
which removed a seven-hour shift of the online lines; EUR conversion; a member bridge (one member per transaction, staff and business
accounts excluded); SKU $\to$ model mapping; national price intervals per SKU; a sales calendar inferred from peaks of the
promotion share plus public holidays (SG 53, TH 106, VN 58 days); and identity stitching that links web cookies to members through
web orders and logins, valid within $\pm 90$ days of the link. 5.58\,M netted lines are refused with one of six reasons (mostly
non-merchandise and cancelled orders) and kept in quarantine.

\textit{Gold respects time.} Eight Gold models build the anchor grid (Mondays after a member's second purchase), three feature
entities, member at $t$ (recency, frequency, monetary value, price profile, web activity), model $\times$ country at $t$ (price,
promotion, season, popularity; 4.29\,M rows) and pair $(u, m, t)$ (own and department history, co-purchase, web intent), the rule
candidates R1--R5, and the labels: every non-returned model bought in $(t, t+\Delta]$ with $\Delta = 30$ days ($\Delta = 90$ as an
auxiliary label). Prices at an anchor come from an ASOF join on the SKU intervals. Every feature reads only data recorded before
its anchor, and a web event counts only if its member link was known at the anchor.

\textit{Quality gates.} Seven checks run across the tiers: null keys (QC-1), referential integrity (QC-2), value domains (QC-3),
uniqueness (QC-4), volume drift (QC-5), point in time (QC-6: six singular tests and a leakage ledger in which every violating row is
logged) and batch--stream reconciliation for an optional streaming lane (QC-7). The 31 dbt models carry 182 tests in five layers.

\subsection{After Processing}
The fact table keeps 66.1\,M net sale lines of the 89.3\,M raw lines, in 26.3\,M transactions by 4.7\,M buying members. Stitching
raises the share of web events linked to a member by 6--8 points (SG 36 to 42\,\%, TH 34 to 40\,\%, VN 30 to 38\,\%). The 64 Gold
tests pass and the QC-6 ledger holds zero point-in-time violations in all three countries. Per country and year, go / no-go gates
on member link ($\ge 60\,\%$), item match ($\ge 95\,\%$), price coverage ($\ge 90\,\%$) and members with two or more transactions
($\ge 20\,000$) group the countries (Table~\ref{tab:gates}): TH and VN pass every year and form the main evaluation set; SG is
processed in full but flagged REVIEW, because its member link is low and its labels miss the repeat purchases of unlinked
transactions. Unless stated otherwise, every number below is on the test split, TH + VN, $\Delta = 30$.

\begin{table}[t]
\caption{Go / no-go gates per country and year (QC-2 to QC-4) and the share of transactions linked to a member.}
\label{tab:gates}
\setlength{\tabcolsep}{4pt}
\centering
\begin{adjustbox}{max width=\columnwidth}
\begin{tabular}{lccccr}
\toprule
Country & 2023 & 2024 & 2025 & 2026 & member link \\
\midrule
TH & PASS   & PASS   & PASS & PASS   & 86--100\,\% \\
VN & PASS   & PASS   & PASS & PASS   & 62--76\,\% \\
SG & REVIEW & REVIEW & PASS & REVIEW & 50--64\,\% \\
\bottomrule
\end{tabular}
\end{adjustbox}
\end{table}

\subsection{Model Inputs and Splits}
The unit of prediction is a (member, anchor Monday) pair. Splits are by anchor date with a 42-day purge gap: train up to
2025-09-15, valid anchors 2025-10-27 to 2025-12-22, test anchors 2026-02-02 to 2026-07-27. Anchors are capped per country by hash
rank (50\,000 train, 15\,000 valid, 50\,000 test).

Every model-generated table carries its fit cutoff, the latest source timestamp it read and a model version, and is checked by QC-6:
no row may read an event at or after its anchor, and no fit may use a cutoff later than the anchor. Learned candidate sources and
signal models are refitted at 13 quarterly cutoffs (2023-07-01 to 2026-07-01) in every split; an anchor uses the latest cutoff
$\le t$. A single fit at the end of train would leave 27--34\,\% of the test positives outside a 5\,000-model universe that is
5--10 months old; quarterly refits halve that share, to 15.5\,\% (TH + VN, Section~\ref{sec:retrieval}). A fit at a cutoff inside
valid or test reads the purchases of that split before the cutoff, so the 42-day purge gap separates the splits for the rankers but
not for the candidate sources and signals; their protection is the point-in-time rule above. Each country has its own sources,
universe and history: a member's purchases in another country are not read. Rankers are trained once on train and selected on
valid. Fig.~\ref{fig:arch} shows the two-stage system that the following sections build on these inputs, and Fig.~\ref{fig:pit_timeline}
contrasts the rolling quarterly refit schedule of the retrieval sources and stacked signals against the single training-split fit of the rankers.

\Figure[t!](topskip=0pt, botskip=0pt, midskip=0pt)[width=0.99\columnwidth]{fig_arch.png}{Two-stage system. The candidate sources and signal models are refitted at quarterly cutoffs before the anchors they score; the rankers are fitted once on train.\label{fig:arch}}

\Figure[t!](topskip=0pt, botskip=0pt, midskip=0pt)[width=0.95\textwidth]{fig_pit_timeline.png}{Temporal split protocol: 42-day purge gaps between anchor splits, 13 quarterly rolling cutoffs for learned retrieval sources (R6--R13) and P2 stacked signals, and a single training-split fit for the rankers (P1, P2, P3), enforced by QC-6 point-in-time assertions.\label{fig:pit_timeline}}

The lakehouse (DuckDB 2.0, dbt 1.12) runs on a single machine: one country on a laptop, the three countries on a 16-thread server
with the same commands. The models run on one shared GH200 node (72 Arm cores and one Hopper-class GPU); deep models use bfloat16
mixed precision. Every model with hyper-parameters receives
the same budget of 10 Optuna~\cite{akiba2019optuna} TPE~\cite{bergstra2011tpe} trials on valid, which is TPE's random start-up phase,
i.e., a seeded random search; the rankers were re-tuned with 30 trials as a check (Section~\ref{sec:results}) and the final fits
repeated with three seeds.

% ---------------------------------------------------------------------------------------------------------------------------------
\section{Unconditional Evaluation and Error Analysis}
\label{sec:protocol}

With the inputs fixed, the first half of RQ2 asks where the rule-based baseline fails when every purchase counts.

\subsection{Protocol}
The evaluated unit is every used anchor $a$ with at least one non-returned purchase in $(t, t+\Delta]$; its positives $P_a$ are all of
those purchases, including those outside the candidate list $C_a$, which are placed below every candidate, and those outside the
per-anchor model grid that filters every candidate list (0.04\,\% of the positives). With $L_a^K$ the top $K$ of
the ranked list,
\begin{equation}
\mathrm{Recall}@K = \frac{|P_a \cap L_a^K|}{|P_a|},
\end{equation}
\begin{equation}
\mathrm{NDCG}@K = \frac{\sum_{k=1}^{K} \mathbf{1}[\ell_k \in P_a] / \log_2(k+1)}{\sum_{k=1}^{\min(K, |P_a|)} 1/\log_2(k+1)},
\end{equation}
where $\ell_k$ is the item at rank $k$, so the ideal gain counts every positive, retrieved or not~\cite{jarvelin2002ndcg}. A positive is \emph{repeat} if the member bought
the model on or before $t$ and \emph{explore} otherwise~\cite{li2023nbr} (the features read strictly before $t$; a purchase on the
anchor Monday itself makes 0.3\,\% of the repeats); recall\_repeat@20 and recall\_explore@20 average over the anchors
that have positives of that kind. We also report HR@10 (at least one hit in the top 10) and the candidate recall (share of positives
anywhere in $C_a$). Metrics are member means; confidence intervals come from a bootstrap over members~\cite{efron1993bootstrap}
(500 resamples) and model differences from a paired bootstrap over the same members (1\,000 resamples).

\subsection{The Rule-Based Baseline Under the New Protocol}
Table~\ref{tab:t28} puts both protocols side by side. The conditional protocol drops the anchors whose list holds no purchase
(\nNoPosShare{}\,\% of the anchors with a purchase over the three countries, 38.8\,\% in TH + VN), which doubles the headline number.

\begin{table}[t]
\caption{P1-frozen under the conditional and unconditional protocols (test, $\Delta = 30$).}
\label{tab:t28}
\setlength{\tabcolsep}{3pt}
\centering
\begin{adjustbox}{max width=\columnwidth}
\begin{tabular}{lrrrrrr}
\toprule
Protocol, scope & anchors & NDCG@20 & HR@10 & rec.\ rep. & rec.\ exp. & cand. \\
\midrule
conditional, TH + VN   & 25\,538 & 0.206 & 0.368 & 0.826 & 0.218 & 1.000 \\
unconditional, TH + VN & 41\,734 & \textbf{0.095} & 0.225 & 0.824 & 0.073 & 0.370 \\
unconditional, all     & 63\,574 & 0.090 & 0.217 & 0.797 & 0.065 & 0.351 \\
\bottomrule
\end{tabular}
\end{adjustbox}
\end{table}

\subsection{Where the Errors Are}
Each of the 213\,289 test positives is put in one cell: ranked in the top 10, ranked below 10, proposed by a source but cut by the
200 cap, or proposed by no source (Table~\ref{tab:t28b}, Fig.~\ref{fig:t28b}).

\begin{table}[t]
\caption{Fate of the test positives of P1-frozen over R1--R5 candidates (shares of positives).}
\label{tab:t28b}
\setlength{\tabcolsep}{4pt}
\centering
\begin{adjustbox}{max width=\columnwidth}
\begin{tabular}{lrrrrr}
\toprule
 & positives & top 10 & rank 11+ & capped & no source \\
\midrule
repeat  & 18\,251  & 0.654 & 0.332 & 0.013 & 0.002 \\
explore & 195\,038 & 0.028 & 0.227 & 0.111 & 0.634 \\
all     & 213\,289 & 0.082 & 0.236 & 0.103 & 0.580 \\
\bottomrule
\end{tabular}
\end{adjustbox}
\end{table}

\Figure[t!](topskip=0pt, botskip=0pt, midskip=0pt)[width=0.99\columnwidth]{fig_t28b_errors.png}{Where the test positives of the rule-based baseline end up.\label{fig:t28b}}

Repeat is largely solved; explore is 91\,\% of the positives and 74.5\,\% of it never reaches the ranker. Of the capped positives,
95\,\% come from R5 (the national top 300), whose lower ranks do not fit the cap. The positives without a source are mid-tail, not
long-tail: 95\,\% have a popularity rank between 301 and 5\,000 at the anchor. A learned source over a per-country universe of the
top 5\,000 models can reach most of them, but the universe is fixed at the quarterly cutoff, not at the anchor: it holds 84.5\,\% of
the test positives (TH + VN, 83.4\,\% of the explore ones), which is the ceiling of every learned source (Table~\ref{tab:t28c}). R4 (sport $\times$ offer, recall 0.0002) is dropped:
only 1.07\,\% of the anchors have a favourite sport. These facts set the agenda of Section~\ref{sec:retrieval}.

% ---------------------------------------------------------------------------------------------------------------------------------
\section{Learned Retrieval (RQ2)}
\label{sec:retrieval}

\subsection{Sources}
R1 (repurchase) already holds 99.7\,\% of the repeat positives, so every learned source is \emph{explore-only}: it masks the models the
member bought before the anchor and is tuned on the explore coverage it \emph{adds} to R1--R5@200 on valid ($\Delta$ cand\_explore of
the union), not on its own recall. Each source is fitted per country at each quarterly cutoff over the top 5\,000 models by buyers in
the previous 365 days, and proposes 100 candidates per anchor (Table~\ref{tab:sources}). With 5\,000 models per country, an exact
matrix product on the GPU replaces an approximate nearest-neighbour index.

The two neural sources are adapted to the target, and each of the eight sources represents a distinct retrieval mechanism over the
5\,000-model cutoff universe:
\begin{itemize}
\item \textbf{Linear item--item and windowed regression (R6 EASE, R11 explin).} EASE~\cite{steck2019ease} builds a binary member $\times$
universe matrix $X$ over the 730 days before the cutoff and solves the closed-form auto-encoder $P = (X^\top X + \lambda I)^{-1}$,
$B = -P / \mathrm{diag}(P)$ with $\mathrm{diag}(B) = 0$, scoring a member's pre-anchor binary history via $s = xB$. Whereas EASE
reconstructs past co-purchases, R11 (\emph{explin}) fits a supervised ridge regression $W = (X^\top X + \lambda I)^{-1} X^\top Y$ directly
from time-decayed history before reference days $t$ ($0.5^{\Delta d / h}$ on each model's latest purchase) to the binary vector $Y$ of
explore purchases in $(t, t+\Delta]$ (up to three reference days per member and $2\cdot 10^6$ rows, with $t+\Delta < c$).
\item \textbf{Sequence embeddings and graph walks (R7 item2vec, R12 RP3beta).} R7 trains skip-gram with negative sampling~\cite{barkan2016item2vec}
over each member's chronological purchase sequence before the cutoff ($d=64$, window 10, 10 negatives, 5 epochs), represents the member at $t$
by the exponentially decayed mean ($0.5^{\Delta d / h}$) of her purchased item vectors, and retrieves by exact cosine similarity. R12
(\emph{RP3beta}~\cite{paudel2017rp3beta}) computes a three-step bipartite graph transition $W_{ij} = \mathrm{deg}(i)^{-\alpha}\bigl(\sum_u x_{ui}x_{uj}\mathrm{deg}(u)^{-\alpha}\bigr)\mathrm{deg}(j)^{-\beta}$
with $\mathrm{diag}(W)=0$, penalising popular hubs via $\beta$.
\item \textbf{Time-decayed nearest neighbours (R8 TIFU-KNN).} Following~\cite{hu2020tifu}, a member's purchase days before $t$ are ordered
newest-first into groups of $m=7$; basket $b$ receives weight $r_g^{\lfloor b/m \rfloor} r_b^{b \bmod m} / m$ to form personal vector $u$,
which is blended as $\alpha u + (1-\alpha)\bar{u}_{\mathcal{N}_k(u)}$ with her $k$ nearest cosine neighbours among up to 300\,000 active members at the cutoff.
\item \textbf{Two-tower and causal basket Transformer (R9, R10).} The two-tower model (R9) encodes the last 50 purchase tokens (models outside the universe mapped to an OOV token) with a bidirectional Transformer (2 blocks, 2 heads), concatenates masked mean pooling with log-scaled RFM counts and a calendar-month embedding, and scores normalized item representations (shared model, sport-department and price-decile embeddings) scaled by a learned temperature $\tau$; up to four reference days per member are resampled each epoch and trained with an exact full-universe softmax that masks already-bought items. SASRec (R10) runs a causal Transformer over the last 50 daily baskets (up to 16 model embeddings averaged per day) plus a $\log_2$ inter-basket gap embedding, trained with a dense all-action explore loss~\cite{pancha2022pinnerformer} where every historical basket predicts all explore purchases in the following $\Delta$ days whose window closes before the cutoff.
\item \textbf{Context-aware trending (R13).} To replace the capped tail of R5, R13 scores each universe model at anchor $t$ by
$\log(1 + \text{buyer-days in } [t-w, t)) + w_{\mathrm{dept}}\log(\text{member dept share} + 0.01) + w_{\mathrm{price}}\mathbf{1}[|\text{item price decile} - \text{member mean decile}| \le 1]$, reading counts strictly before $t$.
\end{itemize}

\begin{table}[t]
\caption{Learned candidate sources and their valid gain in explore coverage over R1--R5@200: best of the first 10 trials, the
parameters used. EASE, TIFU-KNN, explin, RP3beta and trend ran 30 trials; the best of 30 is at most 0.001 higher. item2vec keeps
its defaults (gain of the build).}
\label{tab:sources}
\setlength{\tabcolsep}{3pt}
\centering
\begin{adjustbox}{max width=\columnwidth}
\begin{tabular}{llr}
\toprule
Source & Model (tuned parameters) & $\Delta$ cand.\ exp. \\
\midrule
R6  & EASE, closed-form item--item ($\lambda$)                   & +0.085 \\
R7  & item2vec over purchase sequences (defaults)                & +0.074 \\
R8  & TIFU-KNN ($k$, $r_b$, $r_g$, $\alpha$)                      & +0.057 \\
R9  & two-tower, full-universe masked softmax ($d$, $\tau$, lr, dropout) & +0.074 \\
R10 & SASRec, dense all-action loss ($d$, blocks, dropout, lr)    & +0.085 \\
R11 & explin, windowed linear ($\lambda$, half-life)              & +0.074 \\
R12 & RP3beta, random walk ($\alpha$, $\beta$)                    & +0.088 \\
R13 & context-aware trend (window, weights)                       & +0.069 \\
\bottomrule
\end{tabular}
\end{adjustbox}
\end{table}

\subsection{Mixer and the Cap}
The ranker receives at most $K$ candidates. The mixer interleaves the sources by weighted rank quotas: candidate $i$ gets the key
\begin{equation}
\kappa(i) = \min_{s \,:\, i \in C_s} \frac{r_s(i) - 0.5}{w_s},
\end{equation}
where $r_s(i)$ is its rank in source $s$ and $w_s$ the source's stand-alone valid Recall@100, and the $K$ smallest keys are kept.
The weight is a hits-per-slot heuristic (all positives, three countries), not a parameter tuned for the ranking objective: it gives
R1 0.106, below every learned source (0.13--0.20) and R5 (0.21), although repeats carry most of the NDCG (Table~\ref{tab:t31}); Section~\ref{sec:results} tests a
repeat-first variant. A per-member variant scales $w_{R1}$ by the member's repeat share; it ties with the fixed policy on valid (TH + VN coverage at 300:
0.4841 against 0.4847), so the fixed policy is used. $K = \nK{}$ is chosen from the recall--cost curve (Fig.~\ref{fig:curve}).

\subsection{Coverage}

\begin{table}[t]
\caption{Candidate coverage on test (TH + VN, $\Delta = 30$).}
\label{tab:t30}
\setlength{\tabcolsep}{3pt}
\centering
\begin{adjustbox}{max width=\columnwidth}
\begin{tabular}{lrrrr}
\toprule
Candidate set & list & cand. & repeat & explore \\
\midrule
R1--R5 (reference)                       & 200     & 0.370 & 0.997 & 0.307 \\
R1--R5 $\cup$ R12 RP3beta (best single)  & 277     & 0.439 & 0.997 & 0.386 \\
R1--R5 $\cup$ R6 / R10 / R7              & 260--282 & 0.430--0.435 & 0.997 & 0.375--0.381 \\
R1--R5 $\cup$ all learned (uncapped)     & 483     & 0.541 & 0.997 & 0.499 \\
\textbf{candidates\_v2, mixer, $K = 300$} & 300    & \textbf{0.450} & 0.972 & \textbf{0.400} \\
mixer, $K = 500$                          & 470    & 0.535 & 0.994 & 0.492 \\
\bottomrule
\end{tabular}
\end{adjustbox}
\end{table}

\Figure[t!](topskip=0pt, botskip=0pt, midskip=0pt)[width=0.8\textwidth]{fig_t30_curve.png}{Candidate coverage against the cap $K$; the fixed and per-member mixers overlap at every $K$.\label{fig:curve}}

Coverage rises from \nCandOld{} to \nCandNew{} at $K = \nK{}$ (Table~\ref{tab:t30}), meeting the plan's target of 0.45; explore coverage
rises 1.30 times (\nExpOld{} to \nExpNew{}), short of the target of 2. The cap costs 2.5 points of repeat coverage: R1 is short for most
anchors (median 12 models, 90th percentile 41), but the tail of long R1 lists loses its slots to explore candidates. The sources are complementary: the union of all eight adds 0.17 over R1--R5 while the
best single source adds 0.07. The neural sources do not beat the cheap ones: SASRec ties with EASE, the two-tower model is below,
and RP3beta, the best single source, trains in minutes against about six GPU-hours of tuning and building for each neural source
(Section~\ref{sec:feasibility}). This ranking is partly a budget effect: the neural sources train for a fixed 8 epochs without early
stopping, and their training loss still falls by about 0.2\,\% per epoch at the end. Refitting them on TH at the valid cutoff with 4,
8 and 16 epochs gives $\Delta$ cand.\ exp.\ 0.074 / 0.077 / 0.082 for the two-tower model and 0.077 / 0.084 / 0.090 for SASRec: both
are under-trained at 8 epochs, and twice the epochs add 0.006 of explore coverage (+8\,\%), more than the gap between SASRec and
RP3beta in Table~\ref{tab:sources}.
The candidates and rankers of this paper keep the 8-epoch sources. QC-6 on candidates\_v2 finds no row that reads an event at or after its anchor and no fit cutoff
after the anchor. The pairs of candidates\_v2 are rendered from the same dbt SQL as Gold and match it on a parity day up to
floating-point noise.

\begin{table}[t]
\caption{Fate of the test positives with candidates\_v2 ($K = \nK{}$) and P3 one head (TH + VN, shares of positives). ``No source''
is split by whether the model is inside the 5\,000-model universe of the anchor's cutoff.}
\label{tab:t28c}
\setlength{\tabcolsep}{3pt}
\centering
\begin{adjustbox}{max width=\columnwidth}
\begin{tabular}{lrrrrrr}
\toprule
 & positives & top 20 & rank 21+ & capped & \multicolumn{2}{c}{no source} \\
 & & & & & in univ. & outside \\
\midrule
repeat  & 11\,533  & 0.785 & 0.151 & 0.043 & 0.019 & 0.002 \\
explore & 126\,810 & 0.073 & 0.293 & 0.099 & 0.383 & 0.153 \\
all     & 138\,343 & 0.132 & 0.281 & 0.094 & 0.352 & 0.140 \\
\bottomrule
\end{tabular}
\end{adjustbox}
\end{table}

\textit{What is left.} Table~\ref{tab:t28c} redoes the error analysis of Table~\ref{tab:t28b} on the new lists. The cap is no longer
the main loss: 9.4\,\% of the positives are proposed and cut, and R5 proposes 42\,\% of those (it was 95\,\% with R1--R5). The
ceiling is the universe: it holds 84.5\,\% of the positives (85.9\,\% as a member mean), and 86.0\,\% are reachable at all (inside
it or proposed by a rule source), against a candidate recall of 0.450. Of the remaining gap, 35\,\% of the positives are inside the
universe but proposed by no source, a retrieval problem, and 14\,\% are outside it, which no learned source over this universe can
fix. The 4.3\,\% of repeats that are capped come mostly from long R1 lists cut by the quotas (no learned source proposes
them: they are masked as explore-only).

\subsection{From Recall to NDCG}
With the ranker frozen (P1-frozen's hyper-parameters and original feature set, refitted on the v2 pairs), the retrieval gain reaches
NDCG@20 only partly. The paired difference v2 $-$ v1 is $+0.0025$ $[+0.0018, +0.0032]$ in NDCG@20 and $+0.0052$ in HR@10, with
$+0.0101$ recall\_explore@20 and $-0.0315$ recall\_repeat@20. The gain is significant but small (0.0947 to 0.0973): the frozen ranker
has no feature that tells it which learned source proposed a candidate, it trades repeat recall for explore recall, and most new
explore candidates land below rank 20.

% ---------------------------------------------------------------------------------------------------------------------------------
\section{Ranking Models}
\label{sec:ranking}

All two-stage rankers score the same candidates\_v2 lists ($K = \nK{}$) and are selected on unconditional valid NDCG@20 with 10 TPE trials;
the selection objective averages the three countries (SG included), while the main tables report TH + VN. Training is listwise and so uses only
the train anchors with at least one positive in their candidate list: 44\,795 of the 150\,000 sampled anchors (30\,\%, 13.2\,M pairs). For GBDT trials, LightGBM early stopping (100 rounds, up to 6\,000 trees) monitors conditional valid NDCG@20 on the positive validation groups, whereas Optuna trial selection evaluates the resulting predictions under the unconditional protocol over all validation buyer anchors. The rankers are thus
trained on the conditional population and selected and evaluated on the unconditional one (Fig.~\ref{fig:rankers_compare}).
\begin{itemize}
\item \textbf{P1-frozen}: the baseline LambdaRank of Section~\ref{sec:protocol}, with its hyper-parameters (tuned once on R1--R5
candidates) and its 73 Gold pair features kept fixed, refitted on the v2 pairs; a reference row, not a tuned competitor.
\item \textbf{P1-tuned}: LightGBM LambdaRank ($\text{truncation level } 20$) over every pair column of candidates\_v2 (85 features: the 73 Gold baseline features plus 12 v2 retrieval columns, namely the per-source ranks of R6--R13 and the updated mixer metadata, while excluding identifier, label, anchor-grid bookkeeping and point-in-time audit columns; \texttt{country\_code} and \texttt{promo\_kind} are declared as native categorical features).
\item \textbf{P2 (stacked GBDT)}: P1-tuned augmented with five out-of-time signal columns (\texttt{hazard\_p}, \texttt{wtp\_q50}, \texttt{wtp\_q90}, \texttt{wtp\_gap50}, \texttt{wtp\_gap90}) produced by three LightGBM models refitted under the 13-cutoff rolling protocol: (i) a binary department-hazard model over member $\times$ department rows (departments bought before $t$ plus the country top 10 over the prior 90 days) predicting $P(\text{buy in department } d \text{ during } (t, t+\Delta])$ from member RFM, department recency/cycle/overdue ratio and 30/90-day department popularity trends; and (ii) two LightGBM quantile regressions of willingness to pay fitted on positive anchors, where $\alpha = 0.5$ predicts the median log unit price (\texttt{y\_lp\_q50}) of the member's purchases in $(t, t+\Delta]$ and $\alpha = 0.9$ predicts the \emph{maximum} log unit price (\texttt{y\_lp\_max}), together with the signed differences between the candidate's log local price and each predicted quantile.
\item \textbf{P3 (UMTR)}: a User Multi-Task Ranker (3.73\,M registered parameters, Fig.~\ref{fig:p3_arch}) combining a bidirectional Transformer encoder (2 blocks, 2 heads, Pre-LN, GELU) over the member's last $L=50$ purchase lines (\texttt{member} $\times$ \texttt{model} $\times$ \texttt{day}, not grouped into baskets; shared model embedding across the full sold-model vocabulary, $\log_2$ anchor-recency gap embedding in 17 buckets, and reverse-position embedding counting back from the latest purchase), a context MLP, an item tower over the candidate's embeddings and pair features, and scaled dot-product DIN-style target attention~\cite{zhou2018din} of each candidate over the sequence states.
\item \textbf{Baselines}: five full-catalogue NBR baseline families (six rows in Table~\ref{tab:t31}) that produce their own ranked lists rather than re-ranking candidates\_v2: P-TopFreq, GP-TopFreq~\cite{li2023nbr}, EASE and TIFU-KNN tuned as full-catalogue NBR rankers with their own
10-trial budget on valid NDCG@20, and SASRec reusing the R10 network and hyper-parameters (without a separate NDCG@20 tuning search) in two forms: the explore-trained R10 model with the seen-item mask removed at scoring, and \emph{SASRecAll}, the same network retrained on all window targets with repeats included.
\end{itemize}

\Figure[t!](topskip=0pt, botskip=0pt, midskip=0pt)[width=0.95\textwidth]{fig_rankers_compare.png}{Comparison of the three two-stage rankers (P1-tuned, P2 stacked GBDT, and P3 UMTR) over the shared \texttt{candidates\_v2} pool ($K=300$).\label{fig:rankers_compare}}

\Figure[t!](topskip=0pt, botskip=0pt, midskip=0pt)[width=0.98\textwidth]{fig_p3_arch.png}{Architecture of the deep ranker P3 (UMTR), contrasting the two-head gated formulation (\texttt{p3\_full}) with the single-head listwise formulation (\texttt{p3\_onehead}).\label{fig:p3_arch}}

In detail, P3 transforms the 84 numerical pair columns (all 85 GBDT columns except \texttt{country\_code}, with \texttt{promo\_kind} hashed into $1\,000$ buckets) via $t(v) = \mathrm{sign}(v)\log(1+|v|)$, standardises them using training-split means and standard deviations, maps $\mathrm{NaN} \to 0$, and appends 16 binary missingness flags for columns with $>1\,\%$ training NaNs ($F=100$ dense inputs per candidate). Note an implementation detail of the context branch: \texttt{self.ctx(x[:, 0])} projects the full feature vector of the first candidate slot (whose member-entity columns are identical across the $K$ candidates, but which also includes that slot's candidate columns) rather than slicing member-only columns. The user representation concatenates masked mean sequence pooling, the latest sequence state $h_{L}$, and the context vector $c$ into $u = \mathrm{MLP}_{\mathrm{user}}([\bar{h}, h_L, c]) \in \mathbb{R}^d$. Each candidate $k$ is encoded as $v_k = \mathrm{MLP}_{\mathrm{item}}([e(m_k), d(\mathrm{dept}_k), \mathrm{MLP}_{\mathrm{feat}}(x_k)]) \in \mathbb{R}^d$, attends over the valid sequence states via scaled dot-product attention $a_k = \mathrm{softmax}(v_k h_{1:L}^\top / \sqrt{d})\,h_{1:L}$, and forms the interaction representation $h_k = \mathrm{MLP}_{\mathrm{inter}}([u, v_k, a_k, u \odot v_k])$.

In its two-head mode (\texttt{p3\_full}), P3 has separate linear heads $s^r_i, s^e_i$ for repeat and explore candidates, each normalized by a softmax within its own candidate group ($R_a$: models with \texttt{hist\_buy\_cnt} $> 0$; $E_a$: the others). A member-level gate $g(u) = \sigma(w_g^\top u) \in (0, 1)$ splits the probability mass:
\begin{equation}
p(i) = \begin{cases}
g(u)\, \dfrac{e^{s^{r}_i}}{\sum_{j \in R_a} e^{s^{r}_j}}, & i \in R_a,\\[2ex]
\bigl(1 - g(u)\bigr)\, \dfrac{e^{s^{e}_i}}{\sum_{j \in E_a} e^{s^{e}_j}}, & i \in E_a.
\end{cases}
\end{equation}
When one group is empty (a member with no history has $R_a = \emptyset$; a list may have $E_a = \emptyset$), the effective gate is set to 0 or 1 (clamped to $[10^{-6}, 1-10^{-6}]$)
and the other group takes the whole mass. The training loss is listwise negative log-likelihood over in-list positives plus two auxiliary terms, $\mathcal{L} = -\frac{1}{|P_a \cap C_a|}\sum_{i \in P_a \cap C_a} \log p(i) + \lambda_g\,\mathrm{BCE}(g(u), y^{\mathrm{rep}})
+ \lambda_{90}\,\mathcal{L}_{90}$, where $y^{\mathrm{rep}}$ is the repeat share of the member's unconditional positives in $(t, t+30]$ and $\mathcal{L}_{90}$ is binary cross-entropy on the candidate-level $\Delta = 90$ label. Each training run iterates through all 6 epochs with AdamW ($\text{weight decay } 0.01$, gradient clip $1.0$, batch size 256 anchors, bfloat16) and restores the checkpoint with the highest unconditional valid NDCG@20. The single-head variant (\texttt{p3\_onehead}) replaces the two group softmaxes and gate with a single softmax $p(i) = e^{s_i}/\sum_{j \in C_a} e^{s_j}$ over all $K$ candidates (omitting the gate loss). Hyper-parameters ($d, \text{dropout}, \text{lr}, \lambda_{90}, \lambda_g$) are tuned on \texttt{p3\_full}; the ablation variants (\texttt{onehead}, \texttt{fixedgate}, \texttt{noseq}) reuse those selected hyper-parameters without separate tuning searches.

% ---------------------------------------------------------------------------------------------------------------------------------
\section{Results}
\label{sec:results}

\subsection{RQ3: Model Comparison}

\begin{table*}[t]
\caption{Ranking on test (TH + VN, 41\,734 anchors, $\Delta = 30$, unconditional), with 95\,\% member-bootstrap intervals of NDCG@20.
Selection on valid with 10 trials; final fits with seed 0 (seeds 1 and 2 in Table~\ref{tab:seeds}). P1-frozen (v2, 85 features)
keeps the frozen hyper-parameters and round count but reads every pair column of P1-tuned.}
\label{tab:t31}
\centering
\begin{tabular}{lrcrrrr}
\toprule
Model & NDCG@20 & 95\,\% CI & HR@10 & recall\_repeat@20 & recall\_explore@20 & cand.\ recall \\
\midrule
P-TopFreq                 & 0.0615 & [0.0597, 0.0632] & 0.156 & 0.877 & 0.000 & 0.093 \\
GP-TopFreq                & 0.0752 & [0.0734, 0.0769] & 0.184 & 0.878 & 0.041 & 0.262 \\
EASE NBR                  & 0.0548 & [0.0535, 0.0561] & 0.146 & 0.368 & 0.072 & 0.231 \\
TIFU-KNN NBR              & 0.0788 & [0.0770, 0.0806] & 0.187 & 0.885 & 0.043 & 0.234 \\
SASRec NBR (R10, explore-trained) & 0.0606 & [0.0589, 0.0622] & 0.156 & 0.435 & 0.067 & 0.235 \\
SASRec NBR (all targets)  & 0.0751 & [0.0732, 0.0769] & 0.180 & 0.571 & 0.068 & 0.245 \\
\midrule
P1-frozen (R1--R5 candidates) & 0.0947 & [0.0925, 0.0968] & 0.225 & 0.824 & 0.073 & 0.370 \\
P1-frozen (v2)            & 0.0973 & [0.0949, 0.0994] & 0.230 & 0.793 & 0.083 & 0.450 \\
P1-frozen (v2, 85 features) & 0.1001 & [0.0981, 0.1022] & 0.236 & 0.793 & 0.090 & 0.450 \\
P1-tuned (v2)             & 0.1011 & [0.0990, 0.1033] & 0.239 & 0.817 & 0.092 & 0.450 \\
P2 (v2)                   & 0.1013 & [0.0994, 0.1034] & 0.240 & 0.806 & 0.093 & 0.450 \\
P3 two heads              & 0.1009 & [0.0988, 0.1029] & 0.237 & 0.845 & 0.089 & 0.450 \\
\textbf{P3 one head}      & \textbf{0.1020} & [0.0999, 0.1040] & 0.240 & 0.838 & 0.090 & 0.450 \\
\bottomrule
\end{tabular}
\end{table*}

\begin{table}[t]
\caption{Paired member-bootstrap differences on test (TH + VN), $a - b$ [95\,\% CI].}
\label{tab:paired}
\setlength{\tabcolsep}{3pt}
\centering
\begin{adjustbox}{max width=\columnwidth}
\begin{tabular}{lll}
\toprule
$a - b$ & NDCG@20 & rec.\ explore@20 \\
\midrule
P1-tuned $-$ P1-frozen (v2) & +0.0039 [+0.0032, +0.0045] & +0.0085 \\
\quad new features: 85 feat.\ $-$ P1-frozen (v2) & +0.0028 [+0.0021, +0.0035] & +0.0069 \\
\quad tuning: P1-tuned $-$ 85 feat.          & +0.0011 [+0.0005, +0.0016] & +0.0016 \\
P2 $-$ P1-tuned             & +0.0002 [$-$0.0002, +0.0006] & +0.0007 \\
P3 two heads $-$ P1-tuned   & $-$0.0002 [$-$0.0009, +0.0005] & $-$0.0031 \\
P3 one head $-$ P1-tuned    & +0.0009 [+0.0002, +0.0016] & $-$0.0016 \\
P3 one head $-$ P2          & +0.0007 [+0.0000, +0.0014] & $-$0.0023 \\
P1-tuned $-$ TIFU-KNN NBR   & +0.0224 [+0.0214, +0.0234] & +0.0488 \\
\bottomrule
\end{tabular}
\end{adjustbox}
\end{table}

\Figure[t!](topskip=0pt, botskip=0pt, midskip=0pt)[width=0.99\columnwidth]{fig_t31_ndcg.png}{NDCG@20 of every model on test with 95\,\% member-bootstrap intervals.\label{fig:t31}}

Tables~\ref{tab:t31} and~\ref{tab:paired} and Fig.~\ref{fig:t31} are read as answers rather than as a leaderboard.
\begin{itemize}
\item Every two-stage ranker beats every NBR baseline by a wide margin (P1-tuned $-$ TIFU-KNN: $+0.022$, 28\,\%). The strongest baseline
is TIFU-KNN, then GP-TopFreq; EASE and SASRec as full-catalogue rankers give up the repeats (recall\_repeat 0.37--0.57 against
0.88). For the R10 SASRec this is partly by construction, since it was trained on explore targets only and its explore mask is only
removed at scoring: retrained on every model of the window, SASRec gains $+0.0146$ $[+0.0135, +0.0157]$ NDCG@20 (repeat recall 0.43
to 0.57) and ties with GP-TopFreq, still $-0.0036$ $[-0.0049, -0.0025]$ below TIFU-KNN. The conclusions on the baselines do not change.
\item Features matter more than tuning or the model family. P1-tuned gains $+0.0039$ over P1-frozen on the same pairs, but most of
it comes from the new features, not from tuning: the frozen hyper-parameters with the 85 features already give $+0.0028$ (mostly
explore recall, $+0.007$), and tuning adds $+0.0011$ (mostly repeat recall, $+0.024$). The tuning share is the size of the gap
between P3 one head and P1-tuned ($+0.0009$).
\item The deep model and the GBDT are within a thousandth. P3 with one head is the only model significantly above P1-tuned ($+0.0009$,
about 1\,\%), and only just above P2. The two-head P3 ties with P1-tuned: it keeps more repeats ($+0.029$ recall\_repeat@20) and finds
fewer explore purchases ($-0.003$).
\item Stacking adds nothing measurable (P2 $-$ P1-tuned $+0.0002$, interval across zero), consistent with Section~\ref{sec:h1}: the Gold
pair features already hold department recency, cycles and price gaps.
\end{itemize}
On valid, over the three countries of the tuning objective, the order is the same (P3 one head 0.1049, P3 two heads 0.1037, P2
0.1034, P1-tuned 0.1033).

\textit{Seed variance.} The final fits of P1-tuned and both P3 variants were repeated with seeds 1 and 2 at the same hyper-parameters
(Table~\ref{tab:seeds}, Fig.~\ref{fig:seeds}). The seed spread ($\le 0.0003$) is below the gap between P3 one head and P1-tuned; the
seed-matched paired differences are $+0.0009$ ($[+0.0002, +0.0016]$), $+0.0010$ ($[+0.0003, +0.0017]$) and $+0.0006$ ($[-0.0000, +0.0013]$, positive in point estimate at all three seeds and statistically significant at seeds 0 and 1, while seed 2's lower 95\,\% bound touches zero), and one head beats two heads at every seed ($+0.0011$,
$+0.0018$, $+0.0011$). The two-head model is the least stable: its repeat recall moves from 0.79 to 0.85 across seeds, i.e., the
learned gate lands in different places from one run to the next.

\begin{table}[t]
\caption{Seed variance of the final fits (test, TH + VN, NDCG@20).}
\label{tab:seeds}
\setlength{\tabcolsep}{3pt}
\centering
\begin{adjustbox}{max width=\columnwidth}
\begin{tabular}{lccc}
\toprule
Model & seeds 0 / 1 / 2 & mean $\pm$ sd & rec.\ repeat mean $\pm$ sd \\
\midrule
P1-tuned     & 0.1011 / 0.1013 / 0.1016 & 0.1013 $\pm$ 0.0002 & 0.809 $\pm$ 0.007 \\
P3 two heads & 0.1009 / 0.1005 / 0.1011 & 0.1008 $\pm$ 0.0003 & 0.824 $\pm$ 0.030 \\
P3 one head  & 0.1020 / 0.1022 / 0.1022 & 0.1022 $\pm$ 0.0001 & 0.831 $\pm$ 0.008 \\
\bottomrule
\end{tabular}
\end{adjustbox}
\end{table}

\Figure[t!](topskip=0pt, botskip=0pt, midskip=0pt)[width=0.99\columnwidth]{fig_t31_seeds.png}{NDCG@20 of the final fits across three seeds.\label{fig:seeds}}

\textit{Is the cap binding?} $K = 500$ instead of 300 raises test coverage from 0.450 to 0.535 for 1.6 times the pairs. With the
rankers refitted or re-tuned, valid NDCG@20 (TH + VN, the decision data) does not improve (P1-tuned 0.1088 to 0.1079, P3 one head
0.1096 to 0.1097), so $K = 300$ is kept. On test, reported once, P3 one head gains $+0.0012$ $[+0.0007, +0.0017]$ and P1-tuned
$-0.0003$ $[-0.0008, +0.0001]$; the one-head gain is not supported on valid and is not used for selection. The cap is not what
limits accuracy; ranking the hard explore candidates is.

\textit{Are the mixer weights binding?} Placing the whole R1 list first restores the 2.5 points of repeat coverage the quotas cut
(0.972 to 0.997 on test), and a learned mixer (LightGBM over the source ranks) lifts valid coverage to 0.502; P1-tuned on either
gives the same valid NDCG@20 (0.1087 and 0.1084 against 0.1088) and $+0.0004$ on test, so the fixed weights are kept.

\textit{Is the tuning budget binding?} The best 10-trial GBDT and P3 trials sat near their learning-rate floors, so every ranker
was re-tuned with 30 trials and lower floors. Valid NDCG@20 does not improve (changes of at most 0.0004), so the 10-trial fits are
kept; on test only P2 moves ($+0.0007$ $[+0.0003, +0.0011]$, to 0.1020, level with P3 one head on test but not on valid).

\subsection{By Country and Segment}
Table~\ref{tab:segments} splits NDCG@20 by country and by member segment (transactions in the 365 days before the anchor: new
$\le 1$, mid 2, loyal $\ge 3$; all countries). The ordering holds in every country and segment, and the differences between P1-tuned
and P3 stay below 0.002 everywhere. The largest gain of the whole pipeline over the reference is on new members (0.079 to 0.090,
+14\,\%), whose purchases are almost all explore and who benefit most from the learned sources. SG (REVIEW scope) gains as much as TH
and VN. A web / no-web split is not reported: 2--19\,\% of the anchors have web events and the web features carry about 1\,\% of the
SHAP mass (Section~\ref{sec:shap}).

\begin{table}[t]
\caption{NDCG@20 on test by country and segment.}
\label{tab:segments}
\setlength{\tabcolsep}{3pt}
\centering
\begin{adjustbox}{max width=\columnwidth}
\begin{tabular}{lrrrrrr}
\toprule
Model & SG & TH & VN & new & mid & loyal \\
\midrule
TIFU-KNN NBR    & 0.069 & 0.085 & 0.072 & 0.065 & 0.080 & 0.078 \\
P1-frozen (v1)  & 0.082 & 0.101 & 0.088 & 0.079 & 0.094 & 0.095 \\
P1-tuned (v2)   & 0.089 & 0.108 & 0.093 & 0.089 & 0.101 & 0.099 \\
P3 two heads    & 0.089 & 0.107 & 0.093 & 0.088 & 0.101 & 0.100 \\
P3 one head     & 0.090 & 0.109 & 0.094 & 0.090 & 0.101 & 0.100 \\
\bottomrule
\end{tabular}
\end{adjustbox}
\end{table}

\subsection{RQ3: Repeat and Explore Heads}

\begin{table}[t]
\caption{P3 ablation, paired differences against P3 with two heads and a learned gate (test, TH + VN).}
\label{tab:t32}
\setlength{\tabcolsep}{2.5pt}
\centering
\begin{adjustbox}{max width=\columnwidth}
\begin{tabular}{lrrrr}
\toprule
Variant $-$ two heads & NDCG@20 & HR@10 & rec.\ rep. & rec.\ exp. \\
\midrule
one head              & +0.0011 & +0.0025 & $-$0.0073 & +0.0016 \\
fixed gate            & $-$0.0066 & $-$0.0105 & $-$0.2753 & +0.0159 \\
no sequence           & $-$0.0003 & +0.0001 & $-$0.0240 & +0.0020 \\
\bottomrule
\end{tabular}
\end{adjustbox}
\\[1ex]
{\footnotesize NDCG@20 intervals: one head [+0.0007, +0.0016]; fixed gate [$-$0.0074, $-$0.0058]; no sequence [$-$0.0007, +0.0000].
Fixed gate: $g$ = the member's historical repeat share.}
\end{table}

\Figure[t!](topskip=0pt, botskip=0pt, midskip=0pt)[width=0.8\textwidth]{fig_t32_ablation.png}{P3 ablation on test.\label{fig:t32}}

The explicit split does \emph{not} improve explore (Table~\ref{tab:t32}, Fig.~\ref{fig:t32}): the single head is better on NDCG@20
and on explore recall, at a small cost in repeat recall. Within the two-head architecture, the learned gate is much better than a fixed rule: replacing the learned $g$ by the member's
Bayesian-smoothed historical repeat share loses 0.0066 NDCG@20 and 27.5 points of repeat recall, though the single-head variant achieves the best overall score without any gate at all. A plausible reading is that the historical share
(about 10\,\% of positives) under-weights the repeat group for the anchors where a repeat is due, which the learned gate can see from
the context. The learned gate ranks segments correctly (mean $g$: loyal 0.23, mid 0.17, new 0.14) but is not calibrated: it is about
twice the observed repeat share (0.12, 0.08, 0.06) and correlates only weakly with the per-anchor target (0.11--0.19). In the \texttt{noseq} variant, disabling both the sequence encoder and target attention together (\texttt{sequence=False, din=False}) reduces repeat recall by $0.024$ with no significant change in NDCG@20 ($-0.0003$ $[-0.0007, +0.0000]$): the tabular pair features already summarize most of the member's history, and because both sequence modules are switched off jointly, their individual effects are not separated here.

\subsection{A Single-Stage Alternative (REX)}
\label{sec:rex}
The candidate pool holds 45\,\% of the positives (cand\_recall, Table~\ref{tab:t31}), which caps every ranker. REX
(Repeat--Explore network, \texttt{ml/retrieval/rex.py}) removes the candidate stage: one network per country $\times$ cutoff scores the
whole 5\,000-model universe, with the rolling protocol, the grid filter and the top-100 output of the other single-stage rows. The
logit of model $i$ for a member at day $t$ is the sum of four channels learned in one softmax over the universe:
an explore term $u(t)\cdot v_i$ (a basket transformer over the last 50 purchase days with gap and age embeddings, plus RFM and
month, against a model / department / price-bucket item vector); a popularity MLP on the buyers of $i$ in the 7 / 30 / 90 days
before $t$, its trend and the same 30 days a year earlier; a repeat MLP on the member's own history of $i$ (purchase days, recency,
age, decayed counts, repurchase rate and cycle of $i$ before the cutoff), added only for models she bought; and a web MLP on her
product views and cart events of $i$ in the 30 days before $t$, by the day the attribution is known. Training reference days are
drawn uniformly in the 30 days before a purchase (not the first), with $t + \Delta$ before the cutoff, and the target is every
model bought in $(t, t + \Delta]$, repeats included. $d$, blocks, dropout and learning rate are tuned with the same 10-trial budget on
valid NDCG@20; the components are standard (SASRec-style encoder, two-tower scoring, TIFU-style repeat features), the
contribution is their combination in one model over the full catalogue.

\begin{table}[t]
\caption{REX against the best rows of Table~\ref{tab:t31} and its ablation (test, TH + VN, NDCG@20; three training seeds for
REX and P3 one head; ablations with the parameters of REX, paired differences against REX).}
\label{tab:rex}
\setlength{\tabcolsep}{3pt}
\centering
\begin{adjustbox}{max width=\columnwidth}
\begin{tabular}{lrrrr}
\toprule
Model & NDCG@20 & HR@10 & rec.\ rep. & rec.\ exp. \\
\midrule
TIFU-KNN (single stage)   & 0.0788 & 0.1868 & 0.8847 & 0.0430 \\
P1-tuned                  & 0.1011 & 0.2393 & 0.8165 & 0.0918 \\
P3 one head, 3 seeds      & 0.1021 & 0.2398 & 0.8307 & 0.0910 \\
REX, 3 seeds              & 0.1031 & 0.2427 & 0.8051 & 0.0905 \\
\midrule
REX without web           & $-$0.0027 & & & \\
REX without popularity    & $-$0.0155 & & & \\
REX without repeat        & $-$0.0133 & & & \\
REX without sequence      & $-$0.0121 & & & \\
\bottomrule
\end{tabular}
\end{adjustbox}
\\[1ex]
{\footnotesize Seed means. REX $-$ P3 one head, paired by seed: +0.0012 [+0.0001, +0.0023], +0.0006 [$-$0.0004, +0.0017],
+0.0008 [$-$0.0002, +0.0019]; REX $-$ TIFU-KNN +0.0245 [+0.0234, +0.0255]. Every ablation interval excludes 0.}
\end{table}

A single network matches the two-stage pipeline with 13 sources and a learned ranker (Table~\ref{tab:rex}): REX is
$+0.0010$ NDCG@20 above P3 one head on the seed mean, with an interval above zero for one seed of three, and $+0.0245$ above the best
single-stage baseline. It reaches the pipeline's explore recall without its candidate pool and loses 2.6 points of repeat recall. Every channel is needed:
popularity ($-0.0155$), repeat ($-0.0133$) and the sequence encoder ($-0.0121$) carry most of the score, and web adds $0.0027$, the
only dimension beyond purchases that REX uses. Three extensions were rejected on valid (REX: 0.1054, all countries; seeds 0.1040 to
0.1054): the price and promotion of the model on $t - 1$ with the member's price affinity and her own repeat features (0.1055), her
repeat features alone (0.1058) and 16 epochs instead of 8 (0.1015, overfitting at the tuned dropout of 0.008). REX does not replace the pipeline: it uses
no price, store or member attribute, its scores are harder to explain than the GBDT's, and it needs a GPU to train. It shows that the
candidate cap, not the ranker, bounds the pipeline; Section~\ref{sec:hybrid} combines it with the pipeline.

\subsection{REX With the Pipeline and Across Countries}
\label{sec:hybrid}
REX and the pipeline err on different anchors: a perfect re-order of the REX top 100 would give 0.372 on test, a perfect order of
the repeats alone only 0.114. Three combinations were tried, every choice made on valid over the three countries. \emph{Rank
fusion} (\texttt{ml/eval/fusion.py}) scores a model by $\sum_j w_j / (k + \mathrm{rank}_j)$ over $N$ lists. \emph{REX features}:
REX, refitted per country $\times$ cutoff over the anchors of candidates\_v2 (point-in-time checks clean), adds its log-probability
and rank to every pair, and P1 is re-tuned with the same budget. \emph{Pooled REX}: one network per cutoff over SG, TH and VN, with
a country embedding and within-country price percentiles; each country keeps its own universe mask and popularity counts
($d$ up to 384 in the search; selected 384, one block, dropout 0.084).

\begin{table}[t]
\caption{REX combined with the pipeline (test, TH + VN, NDCG@20; reported once after selection on valid; paired differences
against REX, 95\,\% member bootstrap).}
\label{tab:hybrid}
\setlength{\tabcolsep}{3pt}
\centering
\begin{adjustbox}{max width=\columnwidth}
\begin{tabular}{lrrrl}
\toprule
List & valid & NDCG@20 & HR@10 & $-$ REX \\
\midrule
REX (seed 0)                                   & 0.1054 & 0.1033 & 0.2427 & \\
REX, 3 seeds fused                             & 0.1069 & 0.1050 & --     & +0.0017 [0.0013, 0.0022] \\
Pooled REX                                     & 0.1083 & 0.1076 & 0.2534 & +0.0043 \\
P1 + REX features                              & 0.1083 & 0.1060 & --     & +0.0027 \\
P1 + pooled REX features                       & 0.1100 & 0.1091 & 0.2569 & +0.0058 \\
REX $\oplus$ P3 one head                       & 0.1098 & 0.1066 & 0.2503 & +0.0033 [0.0027, 0.0039] \\
\textbf{Pooled REX $\oplus$ P3 $\times$3 $\oplus$ P1 + REX} & 0.1122 & \textbf{0.1102} & 0.2567 & +0.0069 \\
\bottomrule
\end{tabular}
\end{adjustbox}
\\[1ex]
{\footnotesize $\oplus$: weighted rank fusion; last row weights 0.4 / 0.3 / 0.3, $k = 5$. Differences without an interval are
differences of means.}
\end{table}

Each combination adds a few thousandths, the size of every model step so far (Table~\ref{tab:hybrid}). REX features lift P1 by
$+0.0048$ $[+0.0041, +0.0056]$, pooled REX features by $+0.0080$ $[+0.0072, +0.0088]$, mostly on explore recall. Pooling helps the
smallest country most (VN $+0.0036$, TH $+0.0014$, SG $+0.0005$) and replaces the three-seed ensemble with one training. The best
list, 0.1102, is $+0.0082$ over P3 one head. Rejected on valid: supermodel or nature rankings of REX as a fusion prior (within the
seed spread) and nature and family embeddings in the item tower (0.1047 against 0.1054).

\subsection{Coarser Product Levels}
\label{sec:levels}
Model-level NDCG@20 stays near 0.11 whatever the model, and the label window is not the lever (valid, $\Delta$ of 7 to 90 days
or the next basket: 0.106--0.118). The item unit was therefore changed: \texttt{ml/eval/granularity.py} re-runs D3 with the
distinct groups bought in $(\text{anchor}, \text{anchor} + \Delta]$ as positives, at five catalogue levels from supermodel (all
variants of a product) and product nature (the product type) to department; a model list becomes a group list by the best model
score of each group. REX is also trained directly at a level, with the model-level parameters.

\begin{table}[t]
\caption{NDCG@20 by product level (test, TH + VN; most popular, the best model-level list aggregated to the level, REX trained at
the level).}
\label{tab:levels}
\setlength{\tabcolsep}{4pt}
\centering
\begin{tabular}{lrrr}
\toprule
Level & popular & aggregated & REX at level \\
\midrule
Model          & 0.075 & 0.107 & -- \\
Supermodel     & 0.111 & 0.140 & 0.147 \\
Product nature & 0.297 & 0.323 & \textbf{0.372} \\
Retail family  & 0.270 & 0.298 & -- \\
Sub-department & 0.430 & 0.452 & -- \\
Department     & 0.560 & 0.592 & -- \\
\bottomrule
\end{tabular}
\\[1ex]
{\footnotesize Aggregated: the fusion of REX and P3 one head (0.107 at the model level).}
\end{table}

The 0.2 target is met by predicting the product type a member buys next (Table~\ref{tab:levels}): REX at the nature level reaches
0.372, $+0.0491$ $[+0.0473, +0.0508]$ over the aggregated model list and $+0.0750$ over the most popular natures. At sub-department
and department, popularity alone exceeds 0.4. This is a coarser question, so the model level stays the headline metric.

\subsection{H1: Stacking Leakage}
\label{sec:h1}
The department hazard is fitted three ways: rolling (quarterly cutoffs, the only point-in-time-correct variant), naive (fitted once on
train and applied in-sample) and leaky (fitted on every split). Every fit samples at most 120\,000 anchors, so that the variants
differ in which anchors they see, not in how many: naive samples 80\,\% of the 150\,000 train anchors, leaky 120\,000 of the 345\,000
anchors of the three splits, i.e., about 35\,\% of the train anchors. The pseudo-test is the last train quarter (anchors from
2025-07-01 to the end of train), inside the naive and leaky fits but after the training anchors of the fused ranker below. Its AUC on train / pseudo-test / test is 0.681 /
0.703 / 0.708 for rolling, 0.725 / 0.744 / 0.705 for naive, and 0.710 / 0.717 / 0.723 for leaky. The in-sample fit looks 0.04 AUC
better on the anchors it has seen and loses that lead on test; the leaky fit stays better on test because it has seen those labels.

Fused into a ranker (P1-frozen parameters, fitted on train anchors up to 2025-05-31), the signals give the NDCG@20 of
Table~\ref{tab:t29b}. H1 holds in direction on both candidate sets, more clearly on v2: the naive fusion is ahead of rolling on its own
training anchors ($+0.0026$) and on the pseudo-test ($+0.0029$), and behind it on test ($-0.0004$). The leaky fusion is the only one
ahead on test ($+0.0011$ on v2), because its signals were fitted on the test labels. On train and pseudo-test the leaky fusion is
below the naive one although it also fits in-sample: it has seen the labels of about 35\,\% of those anchors against 80\,\% for naive,
so its in-sample inflation is smaller. The inflation is small because the signals are
weak given the Gold features, and the gaps are of the order of the member-bootstrap intervals ($\pm 0.002$ on test, $\pm 0.005$ on the
5\,337 pseudo-test anchors). An in-sample stacked signal would nevertheless have been reported as a gain that does not exist on test.

\begin{table}[t]
\caption{Stacking leakage: unconditional NDCG@20 of a fused ranker by signal protocol (TH + VN).}
\label{tab:t29b}
\setlength{\tabcolsep}{3pt}
\centering
\begin{adjustbox}{max width=\columnwidth}
\begin{tabular}{lrrrrrr}
\toprule
 & \multicolumn{3}{c}{v1 candidates} & \multicolumn{3}{c}{v2 candidates} \\
\cmidrule(lr){2-4}\cmidrule(lr){5-7}
Signals & train & pseudo & test & train & pseudo & test \\
\midrule
none    & 0.1219 & 0.1072 & 0.0941 & 0.1246 & 0.1097 & 0.0961 \\
rolling & 0.1227 & 0.1090 & 0.0945 & 0.1255 & 0.1097 & 0.0964 \\
naive   & 0.1246 & 0.1095 & 0.0934 & 0.1281 & 0.1126 & 0.0959 \\
leaky   & 0.1232 & 0.1087 & 0.0944 & 0.1256 & 0.1091 & 0.0975 \\
\bottomrule
\end{tabular}
\end{adjustbox}
\end{table}

\subsection{What the GBDT Uses}
\label{sec:shap}
Grouped TreeSHAP~\cite{lundberg2020treeshap} over 3\,000 test anchors (Table~\ref{tab:shap}) provides a descriptive attribution of what the fitted tree model relies on: the learned source ranks are
the second largest feature group of P1-tuned (16.5\,\% of the mean absolute SHAP value), consistent with the experimental finding in Table~\ref{tab:paired} that adding the 12 new v2 features accounts for $+0.0028$ of the $+0.0039$ gain over the frozen 73-feature specification. In P2 the five stacked signals take 2.4\,\% of the SHAP mass.

\begin{table}[t]
\caption{P1-tuned: share of mean $|\mathrm{SHAP}|$ by feature group (test, all 11 groups).}
\label{tab:shap}
\setlength{\tabcolsep}{3pt}
\centering
\begin{adjustbox}{max width=\columnwidth}
\begin{tabular}{lrrr}
\toprule
Feature group & all & repeat cand. & explore cand. \\
\midrule
model: popularity / season / stock      & 0.475 & 0.298 & 0.492 \\
retrieval: learned source ranks         & 0.165 & 0.202 & 0.162 \\
retrieval: R1--R5 flags / scores / rank & 0.085 & 0.140 & 0.079 \\
member: price profile                   & 0.080 & 0.031 & 0.085 \\
model: price / promo                    & 0.045 & 0.044 & 0.045 \\
pair: department recency / cycle        & 0.035 & 0.041 & 0.034 \\
member: activity / RFM                  & 0.028 & 0.048 & 0.026 \\
pair: co-purchase / sport / budget      & 0.027 & 0.022 & 0.028 \\
pair: own history                       & 0.027 & 0.155 & 0.014 \\
country                                 & 0.020 & 0.006 & 0.022 \\
web behaviour                           & 0.013 & 0.013 & 0.013 \\
\bottomrule
\end{tabular}
\end{adjustbox}
\end{table}

% ---------------------------------------------------------------------------------------------------------------------------------
\section{Value, Cost and Demonstrator (RQ4)}
\label{sec:feasibility}

\subsection{Value Over Simple Rules}
The best model on valid, P3 with one head, is compared with three references on test HR@10 (Table~\ref{tab:t38}). ``Random'' draws
10 slots uniformly from the 5\,000-model universe (analytic); ``popular'' is the country top 10 by buyers in the 28 days before the
anchor, the same list for every member; ``buy-again'' is P-TopFreq. The model is 1.8 times the best-sellers and 1.5 times ``buy the
same again'' overall. On repeats it matches buy-again (both about 0.78); on explore purchases it is slightly below the best-sellers
(0.110 against 0.115). Its value is that it does both at once in one list of ten, which no single rule does.

\begin{table}[t]
\caption{HR@10 of P3 one head against business references (test, TH + VN).}
\label{tab:t38}
\setlength{\tabcolsep}{3pt}
\centering
\begin{adjustbox}{max width=\columnwidth}
\begin{tabular}{lrrrrrr}
\toprule
Positives & P3 & random & popular & buy-again & $\times$ pop. & $\times$ again \\
\midrule
all     & 0.240 & 0.007 & 0.137 & 0.156 & 1.76 & 1.54 \\
repeat  & 0.775 & 0.003 & 0.178 & 0.777 & 4.36 & 1.00 \\
explore & 0.110 & 0.006 & 0.115 & 0.000 & 0.96 & --- \\
\bottomrule
\end{tabular}
\end{adjustbox}
\end{table}

As an offline proxy, not a causal estimate, the GMV-weighted Recall@10 is 0.077 against 0.053 for buy-again and 0.023 for popular,
i.e., 268\,EUR of observed next-30-day purchase value per 1\,000 recommendations against 184 and 79. The denominator counts only the
anchors with a purchase in the window (41.7\,\% of the test anchors); per 1\,000 recommendations to every anchor, buyers or not, the
figures are 112, 77 and 33\,EUR. These purchases happened anyway;
an online test is needed to know how many a recommendation would cause.

\textit{Purchase history only.} The status quo recommends from purchase history alone. To separate what the extra data dimensions
bring from what the two-stage pipeline brings, P1 was re-tuned (10 trials) and refitted on the same pairs with only the 39 of 85
features a purchase log yields: R1, R2, R5 and the learned source ranks, own and department history, co-purchase, sales popularity,
season and model age, activity counts and country. Web behaviour (R3 and the mixed candidate rank), prices and promotions, stock,
stores, calendar, sport affinity and the member price profile are removed. Table~\ref{tab:hist} crosses data and method. The
pipeline on history alone gains $+0.0214$ $[+0.0204, +0.0225]$ NDCG@20 over GP-TopFreq; the other dimensions add $+0.0045$
$[+0.0039, +0.0052]$ for P1-tuned and $+0.0054$ $[+0.0046, +0.0062]$ for P3 one head over P1-hist (paired, test; valid gives the
same order: 0.0819, 0.1044, 0.1088, 0.1096). Of the $+0.0268$ of P3 one head over GP-TopFreq, about four fifths are therefore reached
by retrieval and a learned ranker on the purchase log, and one fifth needs the other data, mostly as explore recall (P1-tuned $-$
P1-hist $+0.0077$, repeat $-0.0086$). The history-only row is slightly optimistic: its candidate pool is the multi-dimensional one
(it still holds the web-intent candidates, without the features that explain them), and the two-tower and trend sources use the
unit price paid on the sales lines.

\begin{table}[t]
\caption{Data $\times$ method on test (TH + VN, NDCG@20). P1-hist: P1 re-tuned on the 39 purchase-history features.}
\label{tab:hist}
\setlength{\tabcolsep}{3pt}
\centering
\begin{adjustbox}{max width=\columnwidth}
\begin{tabular}{lllr}
\toprule
Data & Method & Model & NDCG@20 \\
\midrule
purchase history & frequency rule & GP-TopFreq & 0.0752 \\
                 & NBR            & TIFU-KNN   & 0.0788 \\
                 & pipeline       & P1-hist    & 0.0966 \\
\midrule
multi-dimensional & pipeline      & P1-tuned    & 0.1011 \\
                  &               & P3 one head & \textbf{0.1020} \\
\bottomrule
\end{tabular}
\end{adjustbox}
\end{table}

\subsection{Serving Coverage}
Every evaluated test anchor gets a personalized list of 300 candidates (median 300; all anchors have at least 10). This holds by
construction, because the anchor grid starts after a member's second purchase, so it does not answer RQ4. Over the real sales of the
test window, identified members make 93\,\% of the GMV in TH and 81\,\% in VN, and members with a purchase before the sale day, who
get the personalized path, make 71\,\% and 56\,\% (68\,\% and 58\,\% of the buying members). The rest, anonymous sales and first
purchases, gets the country best-sellers: the models of this paper serve about two thirds of the sales value.

\subsection{Cost}
Table~\ref{tab:cost} lists the cost of the candidate sources on the shared GH200 (wall times are upper bounds). The neural sources cost
50--100 times more compute than EASE or RP3beta for the same or a lower coverage gain. Among the rankers the costs are all small at this
scale (13.2\,M training pairs from the 44\,795 train anchors with a positive): tuning takes 0.22--0.38\,h, a final fit 1.3--5.1\,min, and P3 (3.7\,M
parameters) scores 1\,000 test anchors in 0.09\,s on the GPU.

\begin{table}[t]
\caption{Cost of the candidate sources (13 cutoffs $\times$ 3 countries).}
\label{tab:cost}
\setlength{\tabcolsep}{3pt}
\centering
\begin{adjustbox}{max width=\columnwidth}
\begin{tabular}{llrrrr}
\toprule
Source & device & tune h & build h & $\Delta$ exp. & per build h \\
\midrule
R6 EASE       & GPU & 0.04 & 0.07 & +0.085 & 1.27 \\
R12 RP3beta   & GPU & 0.05 & 0.08 & +0.088 & 1.16 \\
R8 TIFU-KNN   & GPU & 0.04 & 0.08 & +0.057 & 0.75 \\
R13 trend     & CPU & 0.03 & 0.13 & +0.069 & 0.54 \\
R11 explin    & GPU & 0.07 & 0.47 & +0.074 & 0.16 \\
R7 item2vec   & CPU & ---  & 0.56 & ---    & --- \\
R9 two-tower  & GPU & 2.48 & 3.05 & +0.074 & 0.02 \\
R10 SASRec    & GPU & 4.19 & 3.18 & +0.085 & 0.03 \\
\bottomrule
\end{tabular}
\end{adjustbox}
\end{table}

\subsection{Demonstrator}
A static page on a public Hugging Face Space runs the recommender in the browser (Pyodide). The user uploads a purchase file
(member, country, model, purchase date, quantity); the page validates it and returns the top $K$ of each member with the source of
each item (repeat with time decay, co-purchase from EASE neighbours, recent country trend), merged with the same rank-quota rule as
the mixer, and a reason in words. Model identifiers are salted hashes, departments carry generic labels, popularity is shown as a
percentile, the sample members are synthetic, and the salt never leaves the server. This part uses reference tables frozen at
2026-07-28 and shows the retrieval layer and its explanations.

The best model needs inputs an uploaded purchase file does not hold: the web, price, promotion, stock and store features of the
Gold layer and the ranks of eight learned sources, two of them (and P3 itself) PyTorch networks that do not run in the browser. It
is therefore served as a real retailer would serve it, by batch scoring in the warehouse. A second tab shows that batch output:
for 500 test anchors drawn at random (250 per country, one per member), the top 10 of P3 one head next to the top 10 of
GP-TopFreq, the member's purchases before the anchor, and the models she bought in the next 30 days. On this sample the share of
members with at least one of the ten bought is 0.240 for P3 against 0.190 for GP-TopFreq (whole test: 0.240 against 0.184). The
sample is published with the data owner's consent and minimised: members get sequential identifiers in shuffled order, models
the salted hashes and departments the generic labels of the reference, dates become days before the anchor, the history is cut to
the last 40 lines, and no price, store, web or member attribute is included. Every deployment of the Space first scans the published files, the batch output included, for
raw model codes and member identifiers, and locally for a list of organisation names.

% ---------------------------------------------------------------------------------------------------------------------------------
\section{Discussion and Limitations}
\label{sec:discussion}

\textit{GBDT against deep models.} Under the same budget and the same candidates, the deep multi-task ranker gains about 1\,\% NDCG@20
over a tuned GBDT and only in its single-head form, in line with the tabular-data literature~\cite{grinsztajn2022tabular}. The effort paid
most in retrieval (+22\,\% relative coverage) and in features that let the ranker read it (+3\,\% NDCG@20 for the learned source ranks
and the other new pair columns under the frozen parameters); tuning added 1\,\%.

\textit{Method against data.} Most of the gain over the status quo comes from the method, not from the extra data: on purchase history
alone the pipeline reaches 0.0966 against 0.0752 for GP-TopFreq, and the web, price, promotion, stock and store features add 0.005
NDCG@20 on top. A retailer with only a sales log can therefore take most of the benefit; the other dimensions matter for the
explore purchases.

\textit{Explore is still the bottleneck.} Explore coverage rose 1.30 times, not 2, and at the top 10 the model is no better than the
best-sellers on explore purchases. On average 60\,\% of a member's explore purchases are still outside the 300 candidates. Raising the cap to 500
recovers coverage (0.535) but not NDCG@20 on valid: the extra candidates are explore models the rankers cannot yet place in the top
20.

\textit{Model against product type.} Every model step, from the tuned rankers to REX and the fusions, moved model-level NDCG@20 by a
few thousandths, up to 0.1102; the same lists aggregated to the product nature score about three times higher, and REX trained at
that level reaches 0.372. A target of 0.2 is thus a choice of question rather than of model: which product type a member buys next
is predictable, which variant within it much less so.

\textit{Tuning budget.} The selected fits use 10 trials, and a 30-trial re-tuning changed valid NDCG@20 by at most 0.0004. With
Optuna's ten random start-up trials and one sampler seed, the 10-trial budget is a seeded random search; the P3 ablation variants
reused the \texttt{p3\_full} parameters and the sources and NBR baselines kept the 10-trial budget. The margins between P1-tuned, P2
and P3 are below 0.001, but P3 one head stays above P1-tuned at all three seeds and both budgets.

\textit{Offline only.} GMV figures are a proxy; recommendations change behaviour, which an offline replay cannot measure.

\textit{Scope.} SG labels are REVIEW; SG results are reported but not used for conclusions. gSASRec, distillation of P3 into a GBDT,
conformal intervals, drift monitoring, and the what-if and parity check of the demonstrator were dropped from the plan.

\textit{Under-trained neural sources.} The two-tower and SASRec sources gain 8\,\% explore coverage from 16 epochs instead of 8
(TH, valid); the comparison of neural and cheap sources is therefore at a fixed and binding epoch budget, and a rebuild with early
stopping on the valid coverage is the first item of future work.

\textit{Future work.} A full-catalogue universe with monthly refits addresses the coverage ceiling of the 5\,000-model universe
(84.5\,\% of the test positives);
content embeddings of the catalogue text and the web click-stream as a second behaviour sequence, possibly with a generative sequential
model such as HSTU~\cite{zhai2024hstu}, target the explore purchases that no co-purchase signal reaches.

% ---------------------------------------------------------------------------------------------------------------------------------
\section{Conclusion}
\label{sec:conclusion}

Eleven operational exports became a point-in-time feature store on one machine: 66.1\,M sale lines kept with every refused row
traceable, 182 tests and zero leakage violations (RQ1). Measured over every purchase, the rule-based recommender built on it was half
as good as the usual protocol says, and most of its misses were purchases it never saw. Treating retrieval as a model, with eight explore-only sources refitted quarterly and mixed under a 300-candidate cap, raised the
share of purchases within reach from \nCandOld{} to \nCandNew{} (RQ2); with tuned rankers that see the source ranks, NDCG@20 rose from
\nRefUncond{} to \nTuned{}--\nBest{} (+7--8\,\%), more than 25\,\% ahead of the best next-basket baseline (\nBaseline{}). Features
gave more than tuning, and tuning more than the model family: a deep two-head ranker with a learned gate does not beat a tuned GBDT,
whereas a single-head version without a gate does by 1\,\% (RQ3). Within the two-head formulation, learning the gate outperforms a fixed historical-share rule, while the explicit repeat/explore head split itself is unnecessary. The best model is 1.8 times the best-sellers and 1.5 times ``buy the same again'' in one list of ten, and reaches
about two thirds of the sales value; on purchase history alone the pipeline keeps about four fifths of its gain over the frequency
rules, and the other data dimensions add the rest (RQ4). A single network over the whole catalogue (REX) matches the pipeline,
and fused with it, pooled over the three countries, gives 0.1102; at the product-type level the same approach reaches 0.372. Naive stacking shows the inflation H1 predicted, small because the signals add little over a good feature store.

\section*{Acknowledgment}
This work was carried out as the project of the course Advanced Machine Learning and Applications (IT2043.26.1.CH.01) of the
Master of Information Technology programme at the University of Information Technology, Vietnam National University Ho Chi Minh City,
under the supervision of Dr. Ha Minh Tan, whom the authors thank for the guidance. The data set was made available for academic use by its owner; no organisation-specific
identifier appears in this report or in the repository.

\bibliographystyle{IEEEtran}
\bibliography{custom}
\EOD

\end{document}
